\documentclass[a4,11pt]{aleph-notas}
% Se puede ver la documentación aquí:
% https://github.com/alephsub0/LaTeX_aleph-notas

% -- Paquetes adicionales
\usepackage{enumitem}
\usepackage{aleph-comandos}
\usepackage{booktabs}
\input{codigo-python.tex}
 


% -- Datos
\institucion{Facultad de Ciencias Exactas, Naturales y Ambientales}
\carrera{Catálogo STEM}
\asignatura{Métodos Numéricos}
\tema{Resumen no. 5: Regresión, diferenciación e integración}
\autor{Mario Cueva - Andrés Merino}
\fecha{Periodo 2026-1}

\logouno[0.18\textwidth]{Logos/logoPUCE_04_ac}
\definecolor{colortext}{HTML}{1957d1}
\definecolor{colordef}{HTML}{1957d1}
\fuente{montserrat}


% -- Comandos adicionales
\setlist[enumerate]{label=\roman*.}


\begin{document}

\encabezado


%%%%%%%%%%%%%%%%%%%%%%%%%%%%%%%%%%%%%%
\section{Regresión lineal y polinomial}

\begin{defi}[Regresión lineal]
    Dado un conjunto de datos $(x_1,y_1),(x_2,y_2),\ldots,(x_n,y_n)$, la \textbf{regresión lineal} busca encontrar el modelo
    \[
        y = a_0 + a_1 x
    \]
    que mejor represente la tendencia de los datos en el sentido de los mínimos cuadrados.
\end{defi}

\begin{defi}[Ajuste por mínimos cuadrados]
    El criterio de \textbf{mínimos cuadrados} determina los coeficientes $a_0$ y $a_1$ que minimizan la suma de los cuadrados de los residuos:
    \[
        S = \sum_{i=1}^{n}(y_i - a_0 - a_1 x_i)^2.
    \]
    Los valores óptimos se obtienen de las ecuaciones normales:
    \[
        a_1 = \frac{n\displaystyle\sum_{i=1}^{n}x_i y_i - \displaystyle\sum_{i=1}^{n}x_i\displaystyle\sum_{i=1}^{n}y_i}{n\displaystyle\sum_{i=1}^{n}x_i^2 - \left(\displaystyle\sum_{i=1}^{n}x_i\right)^2},
        \qquad
        a_0 = \bar{y} - a_1\bar{x},
    \]
    donde $\bar{x}$ y $\bar{y}$ son los promedios de $x_i$ e $y_i$, respectivamente.
\end{defi}

\begin{ejem}
    Consideremos los datos
    \[
        \begin{array}{c|cccc}
            x_i & 1 & 2 & 4 & 5\\
            \hline
            y_i & 2 & 3 & 5 & 7
        \end{array}
    \]
    Se tiene $n=4$, $\sum x_i=12$, $\sum y_i=17$, $\sum x_i^2=46$ y $\sum x_iy_i=63$. Entonces
    \[
        a_1 = \frac{4(63)-12(17)}{4(46)-12^2} = \frac{252-204}{184-144} = \frac{48}{40} = 1.2,
        \qquad
        a_0 = \frac{17}{4}-1.2\cdot\frac{12}{4} = 4.25-3.60 = 0.65.
    \]
    El modelo ajustado es $y = 0.65 + 1.2x$.
\end{ejem}

\begin{defi}[Linealización de datos]
    Cuando la relación entre las variables no es lineal, es posible transformar los datos para obtener un modelo lineal equivalente. Algunas transformaciones comunes son:
    \begin{enumerate}
        \item Modelo exponencial $y=ae^{bx}$: se toma $Y=\ln(y)$, obteniendo $Y=\ln(a)+bx$.
        \item Modelo potencial $y=ax^b$: se toma $Y=\ln(y)$ y $X=\ln(x)$, obteniendo $Y=\ln(a)+bX$.
        \item Modelo recíproco $y=\dfrac{1}{a+bx}$: se toma $Y=\dfrac{1}{y}$, obteniendo $Y=a+bx$.
    \end{enumerate}
    Una vez transformados los datos, se aplica regresión lineal sobre las variables transformadas y se recuperan los parámetros originales.
\end{defi}

\begin{ejem}
    Supongamos que los datos siguen un modelo exponencial $y=ae^{bx}$. Al aplicar la transformación $Y=\ln(y)$, los datos transformados son:
    \[
        \begin{array}{c|ccc}
            x_i & 0 & 1 & 2\\
            \hline
            Y_i=\ln(y_i) & 0.00 & 1.00 & 2.00
        \end{array}
    \]
    La regresión lineal sobre $(x_i,Y_i)$ da $\ln(a)=0$ y $b=1$, por lo que el modelo ajustado es $y=e^x$.
\end{ejem}

\begin{defi}[Regresión polinomial]
    La \textbf{regresión polinomial} de grado $m$ ajusta el modelo
    \[
        y = a_0 + a_1 x + a_2 x^2 + \cdots + a_m x^m
    \]
    a un conjunto de $n>m$ datos. Los coeficientes se obtienen minimizando la suma de cuadrados de los residuos, lo que conduce a un sistema lineal de ecuaciones normales expresable en forma matricial como
    \[
        A^T A\,\mathbf{a} = A^T\mathbf{y},
    \]
    donde $A$ es la matriz de Vandermonde con columnas $1,x,x^2,\ldots,x^m$ y $\mathbf{a}=(a_0,a_1,\ldots,a_m)^T$.
\end{defi}

\begin{ejem}
    Para los datos
    \[
        \begin{array}{c|ccccc}
            x_i & 0 & 1 & 2 & 3 & 4\\
            \hline
            y_i & 2.1 & 3.0 & 6.2 & 11.5 & 18.8
        \end{array}
    \]
    se ajusta el modelo cuadrático $y=a_0+a_1 x+a_2 x^2$. Resolviendo las ecuaciones normales se obtiene $a_0\approx 1.96$, $a_1\approx 0.60$, $a_2\approx 1.12$.
\end{ejem}

\begin{defi}[Regresión lineal múltiple]
    La \textbf{regresión lineal múltiple} ajusta un modelo que depende de $k$ variables independientes $x_1,x_2,\ldots,x_k$:
    \[
        y = a_0 + a_1 x_1 + a_2 x_2 + \cdots + a_k x_k.
    \]
    Los coeficientes se determinan resolviendo el sistema normal $A^T A\,\mathbf{a} = A^T\mathbf{y}$, donde $A$ es la matriz con una columna de unos seguida de $k$ columnas con los valores de cada variable independiente.
\end{defi}

\begin{ejem}
    Para el modelo $z=a_0+a_1 x+a_2 y$ con los datos
    \[
        \begin{array}{c|cccc}
            x_i & 1 & 2 & 3 & 4\\
            \hline
            y_i & 1 & 3 & 2 & 4\\
            \hline
            z_i & 3.1 & 7.8 & 8.2 & 14.0
        \end{array}
    \]
    la matriz de diseño es
    \[
        A=
        \begin{pmatrix}
            1 & 1 & 1\\
            1 & 2 & 3\\
            1 & 3 & 2\\
            1 & 4 & 4
        \end{pmatrix}.
    \]
    Resolviendo $A^T A\,\mathbf{a}=A^T\mathbf{z}$ se obtienen los coeficientes $a_0$, $a_1$ y $a_2$ del modelo.
\end{ejem}

\begin{advertencia}
    La calidad del ajuste puede medirse con el coeficiente de determinación $R^2\in[0,1]$. Un valor cercano a $1$ indica que el modelo explica gran parte de la variabilidad de los datos, aunque esto no garantiza que sea adecuado para predecir fuera del rango observado.
\end{advertencia}


%%%%%%%%%%%%%%%%%%%%%%%%%%%%%%%%%%%%%%
\section{Diferenciación numérica}

\begin{defi}[Diferenciación numérica]
    La \textbf{diferenciación numérica} aproxima la derivada de una función $f$ en un punto $x$ usando valores de $f$ en puntos cercanos. Se basa en el cociente diferencial:
    \[
        f'(x) = \lim_{h\to 0}\frac{f(x+h)-f(x)}{h}.
    \]
    Al sustituir el límite por un valor pequeño pero finito de $h$, se obtienen fórmulas llamadas \textbf{diferencias finitas}.
\end{defi}

\begin{defi}[Diferencia hacia adelante]
    La \textbf{diferencia finita hacia adelante} aproxima la derivada mediante
    \[
        f'(x)\approx\frac{f(x+h)-f(x)}{h}.
    \]
    El error de truncamiento es $O(h)$.
\end{defi}

\begin{ejem}
    Para $f(x)=e^x$ en $x=1$ con $h=0.1$:
    \[
        f'(1)\approx\frac{e^{1.1}-e^{1}}{0.1}=\frac{3.0042-2.7183}{0.1}=2.859.
    \]
    El valor exacto es $f'(1)=e\approx 2.718$, por lo que el error es $|2.718-2.859|\approx 0.141$.
\end{ejem}

\begin{defi}[Diferencia hacia atrás]
    La \textbf{diferencia finita hacia atrás} aproxima la derivada mediante
    \[
        f'(x)\approx\frac{f(x)-f(x-h)}{h}.
    \]
    El error de truncamiento es también $O(h)$.
\end{defi}

\begin{ejem}
    Para $f(x)=e^x$ en $x=1$ con $h=0.1$:
    \[
        f'(1)\approx\frac{e^{1}-e^{0.9}}{0.1}=\frac{2.7183-2.4596}{0.1}=2.587.
    \]
    El error es $|2.718-2.587|\approx 0.131$.
\end{ejem}

\begin{defi}[Diferencia centrada]
    La \textbf{diferencia finita centrada} aproxima la derivada usando puntos a ambos lados de $x$:
    \[
        f'(x)\approx\frac{f(x+h)-f(x-h)}{2h}.
    \]
    El error de truncamiento es $O(h^2)$, mejor que los métodos anteriores.
\end{defi}

\begin{ejem}
    Para $f(x)=e^x$ en $x=1$ con $h=0.1$:
    \[
        f'(1)\approx\frac{e^{1.1}-e^{0.9}}{0.2}=\frac{3.0042-2.4596}{0.2}=2.723.
    \]
    El error es $|2.718-2.723|\approx 0.005$, significativamente menor que los métodos anteriores.
\end{ejem}

Los tres métodos se comparan en la siguiente tabla:

\begin{center}
\small
\begin{tabular}{p{0.26\textwidth}p{0.40\textwidth}p{0.18\textwidth}}
\toprule
Método & Fórmula & Error \\
\midrule
Diferencia hacia adelante & $\dfrac{f(x+h)-f(x)}{h}$ & $O(h)$ \\[4mm]
Diferencia hacia atrás & $\dfrac{f(x)-f(x-h)}{h}$ & $O(h)$ \\[4mm]
Diferencia centrada & $\dfrac{f(x+h)-f(x-h)}{2h}$ & $O(h^2)$ \\[2mm]
\bottomrule
\end{tabular}
\end{center}

\begin{advertencia}
    Un valor de $h$ muy pequeño no siempre mejora la aproximación: los errores de redondeo se amplifican cuando $h\to 0$, por lo que existe un valor óptimo de $h$ que equilibra el error de truncamiento y el error numérico.
\end{advertencia}


%%%%%%%%%%%%%%%%%%%%%%%%%%%%%%%%%%%%%%
\section{Diferenciación numérica de orden superior y parcial}

\begin{defi}[Segunda derivada numérica]
    La \textbf{segunda derivada} de $f$ en $x$ puede aproximarse combinando evaluaciones de $f$. La fórmula centrada es
    \[
        f''(x)\approx\frac{f(x+h)-2f(x)+f(x-h)}{h^2}.
    \]
    El error de truncamiento de esta fórmula es $O(h^2)$.
\end{defi}

\begin{ejem}
    Para $f(x)=e^x$ en $x=1$ con $h=0.1$:
    \[
        f''(1)\approx\frac{e^{1.1}-2e^{1}+e^{0.9}}{(0.1)^2}=\frac{3.0042-5.4366+2.4596}{0.01}=\frac{0.0272}{0.01}=2.72.
    \]
    El valor exacto es $f''(1)=e\approx 2.718$, por lo que el error es $|2.718-2.72|\approx 0.002$.
\end{ejem}

\begin{defi}[Derivada parcial numérica]
    Para una función $f(x,y)$, las \textbf{derivadas parciales} numéricas se obtienen manteniendo una variable fija. Usando diferencias centradas:
    \[
        \frac{\partial f}{\partial x}(x,y)\approx\frac{f(x+h,y)-f(x-h,y)}{2h},
        \qquad
        \frac{\partial f}{\partial y}(x,y)\approx\frac{f(x,y+h)-f(x,y-h)}{2h}.
    \]
\end{defi}

\begin{ejem}
    Para $f(x,y)=x^2+3xy+y^2$ en el punto $(1,2)$ con $h=0.1$:
    \begin{align*}
        f(1.1,2)&=1.21+6.6+4=11.81,\\
        f(0.9,2)&=0.81+5.4+4=10.21.
    \end{align*}
    Entonces
    \[
        \frac{\partial f}{\partial x}(1,2)\approx\frac{11.81-10.21}{2(0.1)}=\frac{1.60}{0.2}=8.0.
    \]
    El valor exacto es $\dfrac{\partial f}{\partial x}=2x+3y=2(1)+3(2)=8$.
\end{ejem}

\begin{advertencia}
    Las mismas consideraciones sobre el tamaño de $h$ aplican para derivadas de orden superior y parciales. Reducir $h$ mejora el error de truncamiento, pero puede amplificar los errores de redondeo.
\end{advertencia}


%%%%%%%%%%%%%%%%%%%%%%%%%%%%%%%%%%%%%%
\newpage
\section{Integración numérica: regla del trapecio}

\begin{defi}[Regla del trapecio simple]
    La \textbf{regla del trapecio simple} aproxima la integral de $f$ en $[a,b]$ por el área del trapecio formado por los puntos extremos:
    \[
        \int_a^b f(x)\,dx \approx \frac{b-a}{2}\bigl[f(a)+f(b)\bigr].
    \]
    El error de esta aproximación es $-\dfrac{(b-a)^3}{12}f''(\xi)$ para algún $\xi\in(a,b)$.
\end{defi}

\begin{ejem}
    Para estimar $\displaystyle\int_0^1 e^x\,dx$ con la regla del trapecio simple:
    \[
        \int_0^1 e^x\,dx\approx\frac{1}{2}\bigl[e^0+e^1\bigr]=\frac{1+e}{2}\approx 1.8591.
    \]
    El valor exacto es $e-1\approx 1.7183$, por lo que el error es $|1.7183-1.8591|\approx 0.141$.
\end{ejem}

\begin{defi}[Regla del trapecio compuesta]
    La \textbf{regla del trapecio compuesta} divide $[a,b]$ en $n$ subintervalos de longitud $h=\dfrac{b-a}{n}$, con nodos $x_i=a+ih$ para $i=0,1,\ldots,n$. La aproximación es
    \[
        \int_a^b f(x)\,dx\approx\frac{h}{2}\left[f(x_0)+2\sum_{i=1}^{n-1}f(x_i)+f(x_n)\right].
    \]
    El error global es $O(h^2)$.
\end{defi}

\begin{ejem}
    Para $\displaystyle\int_0^1 e^x\,dx$ con $n=4$ subintervalos se tiene $h=0.25$ y los nodos son $0,\,0.25,\,0.5,\,0.75,\,1$:
    \[
        \begin{array}{c|ccccc}
            x_i & 0 & 0.25 & 0.5 & 0.75 & 1\\
            \hline
            e^{x_i} & 1.0000 & 1.2840 & 1.6487 & 2.1170 & 2.7183
        \end{array}
    \]
    Entonces
    \[
        \int_0^1 e^x\,dx
        \approx
        \frac{0.25}{2}\bigl[1.0000+2(1.2840+1.6487+2.1170)+2.7183\bigr]
        =
        0.125\times 13.8177
        \approx
        1.7272.
    \]
    El error es $|1.7183-1.7272|\approx 0.009$.
\end{ejem}

\begin{advertencia}
    Al duplicar $n$ (reducir $h$ a la mitad), el error de la regla del trapecio compuesta se reduce aproximadamente a la cuarta parte, pues el error es $O(h^2)$.
\end{advertencia}

\begin{pscodigo}
\begin{function}[H]
\KwData{Una función $f$, extremos $a$ y $b$, número de subintervalos $n$}
\KwResult{Aproximación de $\displaystyle\int_a^b f(x)\,dx$}
\SetKwFunction{trap}{Trapecio}
\Fn{\trap{$f,a,b,n$}}{
    $h:=\dfrac{b-a}{n}$\;
    $S:=f(a)+f(b)$\;
    \For{$i:=1$ \KwTo $n-1$}{
        $S:=S+2\cdot f(a+i\cdot h)$\;
    }
    \Return{$\dfrac{h}{2}\cdot S$}
    }
\end{function}
\end{pscodigo}

\end{document}
